%% kapa: exam=1 round=30 year=2019 period=1 form=A subject=real_estate subject_ko=부동산학원론 seq=32 number=112
%% source: Q-Net 감정평가사 기출문제 (cst00309) artlSeq=5205205 file=extracted/1008/5205205/2019년 제30회 감정평가사 1차시험 문제지/2019년 제30회 감정평가사 1차시험 문제지/2019년도 제30회 감정평가사 1차 1교시 A형.hwp
%% conversion: kordoc 4.24.51 -> tools/gen1.py (auto); review=proofread-sonnet; flags=none
\begin{question}{112}
부동산관리방식에 관한 설명으로 옳지 않은 것은?
\begin{choices}
  \item 자기관리방식은 소유자가 직접 관리하는 방식으로 단독주택이나 소형빌딩과 같은 소규모 부동산에 주로 적용된다.
  \item 위탁관리방식은 부동산관리 전문업체에 위탁해 부동산을 관리하는 방식으로 대형건물의 관리에 유용하다.
  \item 혼합관리방식은 관리 업무 모두를 위탁하지 않고 필요한 부분만 따로 위탁하는 방식이다.
  \item 자기관리방식은 전문성 결여의 가능성이 높으나 신속하고 종합적인 운영관리가 가능하다.
  \item 위탁관리방식은 관리 업무의 전문성과 효율성을 제고할 수 있으며 기밀유지의 장점이 있다.
\end{choices}
\end{question}
\answer{5}
